%
% File: teach/elmagii/formulas/src/formulas.tex [plain TeX code]
% Github: https://github.com/elmagii/formulas/
% Last change: November 8, 2026
%
% This is the primary collection of formulas in the course ``Elektromagnetism
% II, 1TE626 (2022-2026)'' for Engineering Physics, held 2022-2026, at Uppsala
% University, Sweden. This is a compilation of the source present at
% elmagii/formulas/src/formulas.tex, into a compact and handy two-column
% format.
%
% Copyright (C) 2022-2026, Fredrik Jonsson, under Gnu General Public
% License (GPL) v3. See the enclosed LICENSE for details.
%
% This program is free software: you can redistribute it and/or modify
% it under the terms of the GNU General Public License as published by
% the Free Software Foundation, either version 3 of the License, or
% (at your option) any later version.
%
% This program is distributed in the hope that it will be useful,
% but WITHOUT ANY WARRANTY; without even the implied warranty of
% MERCHANTABILITY or FITNESS FOR A PARTICULAR PURPOSE.  See the
% GNU General Public License for more details.
%
% You should have received a copy of the GNU General Public License
% along with this program.  If not, see <https://www.gnu.org/licenses/>.
%
\input macros/epsf.tex
\input macros/eplain.tex
\input amssym % to get the {\Bbb E} font (strikethrough E)
\font\twelvecmbx=cmbx12
\def\iint{\mathop{\int\kern-8pt\int}}
\def\iiint{\mathop{\int\kern-8pt\int\kern-8pt\int}}
\def\oiint{\mathop{\int\kern-8pt\int\kern-13.2pt{\bigcirc}}}
\def\sgn{\mathop{\rm sgn}\nolimits} % sign
\def\Re{\mathop{\rm Re}\nolimits}   % real part
\def\Im{\mathop{\rm Im}\nolimits}   % imaginary part
\def\Tr{\mathop{\rm Tr}\nolimits}   % quantum mechanical trace
\def\cleardoublepage{\ifodd\pageno\vfill\eject~\vskip260pt%
  \centerline{[Page intentionally left blank]} ~\vfill\eject%
  \else\vfill\eject\fi}
\noindent{~\vskip -46pt\hskip-46pt%
  {\epsfxsize=72pt\epsfbox{macros/UU_logo_color.eps}}}
\vskip-46pt\hskip46pt
{\twelvecmbx{Formelsamling}}%
  \writenumberedtocentry{index}{\hskip24pt{Formelsamling}}%
  \nobreak\vskip 0pt\noindent\ignorespaces
%%%%%%%%%%%%%%%%% Begin two-column mode %%%%%%%%%%%%%%%%%%
{\newdimen\fullhsize
\fullhsize=170mm \hsize=82mm
\topskip=0pt\baselineskip=12pt\parskip=0pt\leftskip=0pt\rightskip=0pt
\parindent=0pt
\def\fulline{\hbox to\fullhsize}
\footline={\hfill{\rm\folio}}
\let\lr=L \newbox\leftcolumn
\output={\if L\lr
  \global\setbox\leftcolumn=\columnbox \global\let\lr=R
  \else \doubleformat \global\let\lr=L\fi
  \ifnum\outputpenalty>-20000 \else\dosupereject\fi}
\def\doubleformat{\shipout\vbox{\makeheadline
    \fulline{\box\leftcolumn\hfil\columnbox}
  \makefootline}
  \advancepageno}
\def\columnbox{\leftline{\pagebody}}
\def\eqsection#1{%
  \vskip6pt\goodbreak%
  \hrule width\hsize\vskip1pt%
  \hrule width\hsize
  \vskip4pt%
  \centerline{\bf\vphantom{lq} #1}%
  \vskip4pt%
  \hrule width\hsize\vskip1pt%
  \hrule width\hsize\vskip10pt}
\parskip=0pt
\def\eqtitle#1{\vskip2pt\goodbreak{\it #1}\par}
%\def\trimspacing#1{%
%  \begingroup%
%    \parskip=-10pt%
%%    \displaylines{#1}%
%    #1%
%  \endgroup%
%}

%
% Definition of the vertical padding to use for displayed equations.
% In plain TeX, the syntax <dimension> plus <stretch> minus <shrink> defines
% a "glue" (which is TeX's term for flexible spacing). This "glue" specifies
% three quantities:
%     1. Natural size: the default spacing TeX tries to use.
%     2. Stretch: how much TeX is allowed to increase the spacing.
%     3. Shrink: how much TeX is allowed to decrease the spacing.
%
% For example, "\abovedisplayskip=4pt plus 2pt minus 2pt" means a natural
% spacing of 4 pt, a maximum stretch of 2 pt (up to 6 pt), and a maximum
% shrink of 2 pt (down to 2 pt).
%
% Here, \abovedisplayskip and \belowdisplayskip are used when the preceding
% text line is sufficiently long, while \abovedisplayshortskip and
% \belowdisplayshortskip are used when the preceding line is short enough
% that the equation can be positioned closer to it.
%
% \abovedisplayskip=2pt plus 3pt minus 1pt      % Vertical padding of $$...$$
% \belowdisplayskip=2pt plus 3pt minus 1pt      % Vertical padding of $$...$$
% \abovedisplayshortskip=0pt plus 2pt minus 2pt % Vertical padding of $$...$$
% \belowdisplayshortskip=2pt plus 2pt minus 2pt % Vertical padding of $$...$$
\abovedisplayskip=12pt plus 3pt minus 9pt
\belowdisplayskip=12pt plus 3pt minus 9pt
\abovedisplayshortskip=0pt plus 3pt
\belowdisplayshortskip=7pt plus 3pt minus 4pt

%
% Header of the compiled collection of formulas.
%
\def\versionnumber{1.3}             % Version of this collection of formulas
\def\versiondate{8 november, 2026}  % Timestamp of this collection of formulas
\def\version{v.\versionnumber\ (\versiondate)}
\vskip-4pt
\centerline{Uppsala Universitet}
\centerline{Elektromagnetism II, 1TE626}
\centerline{Sammanst{\"a}lld av Fredrik Jonsson}
\centerline{{\tt github.com/hp35/elmagii}}
\centerline{\version}
\vskip10pt
\input ../formulas/src/formulas.tex
\vfill
\supereject
\if R\lr \null\vfill\eject\fi}
% \cleardoublepage
%%%%%%%%%%%%%% Back to normal single-column %%%%%%%%%%%%%%
